\chapter*{Abstract}
\addcontentsline{toc}{chapter}{Abstract}
Wi-Fi signals fill almost every indoor space, and a person who stands or moves in that space changes how those signals travel from a transmitter to a receiver. The Channel State Information (CSI) reported by a Wi-Fi receiver records this change for every subcarrier of the signal, so an ordinary Wi-Fi link can work as a sensor that needs no camera and no wearable device. Over the last decade, researchers have used CSI to detect motion, to notice people behind walls, to recognise daily activities, to detect falls and to measure breathing. Most of these systems, however, were built on a few specific network cards and routers with several antennas, and many of their reported accuracies were measured in the same environment in which the system was set up or trained.

This research asks how much of this sensing ability can be reached with ordinary commodity Wi-Fi devices that many people already have access to, without any special radio equipment. We review nineteen works on Wi-Fi based presence detection, through-the-wall detection, activity recognition, fall detection and vital sign monitoring, compare their hardware, methods and reported results, and identify the gaps that remain. Based on this review we propose a system with one commodity Wi-Fi transmitter and three commodity Wi-Fi receivers that detects presence and movement, gives a coarse estimate of where the movement is, and is evaluated on recording sessions that were not used for training, so that the reported accuracy is not optimistic.

\vspace{0.5cm}
\noindent\textbf{Keywords:} Wi-Fi Sensing, Channel State Information, Device-Free Sensing, Human Presence Detection, Motion Detection, Commodity Wi-Fi Devices, Internet of Things
